\section{The theorems}
\label{sec:theorems}

Both sides are instances of one structure,
\lean{TransitionSystem} (\file{Core/TransitionSystem.lean}): an initial
state, a partial step function that returns the events of a transition, the
request a label attempts (if any), whether a label is a \emph{progress
step}, and whether a state is \emph{closed}. Finite traces, runs for an
input word, markers, and fair maximal traces are defined once, for any
transition system, so the two trace sets of each statement are built the
same way. Figure~\ref{fig:overview} shows the shape of the claim.

\begin{figure}[t]
\centering
\begin{tikzpicture}[
  >={Stealth[length=5pt]},
  box/.style={draw,rounded corners=2pt,align=center,font=\small,inner sep=5pt,minimum width=3.6cm,minimum height=0.9cm},
  lab/.style={font=\footnotesize,align=center}]
\node[box] (S) at (0,2.6) {\Lang{} program $S$\\[-1pt]{\footnotesize well formed}};
\node[box] (F) at (8.4,2.6) {Dex Flow $\tr{S}$};
\node[box] (TL) at (0,0) {traces over \ObsL{}};
\node[box] (TD) at (8.4,0) {traces over \Obs{}};
\draw[->] (S) -- node[lab,above] {translation $\tr{\cdot}$} node[lab,below] {\S\ref{sec:translation}} (F);
\draw[->] (S) -- node[lab,left] {\Lang{} semantics\\ \S\ref{sec:language}} (TL);
\draw[->] (F) -- node[lab,right] {Dex core semantics\\ under \Eadm{}, \S\ref{sec:dex}} (TD);
\draw[->] (TD) -- node[lab,above] {projection $\pi$} node[lab,below] {equal sets: finite traces;\\ fair maximal traces with \mterm, \mdead, \mdiv} (TL);
\end{tikzpicture}
\caption{The claim. The traces of the translated Flow, projected to
\Lang{}'s alphabet, are the traces of the program.}
\label{fig:overview}
\end{figure}

\subsection{The admissible environment and the projection}

The Dex alphabet has raw identifiers and the theorem's environment is
restricted, so the statement is over a projection, under an admissible
environment.

\paragraph{\Eadm.}
An admissible input word is a finite or infinite sequence of
\emph{attempted} client requests, accepted and rejected alike. It begins
with exactly one request that starts the Flow at the entry Step with the
program input. Every later request publishes one message to one of the
program's external Channels. The start input and every message are encoded
\Lang{} literals (\S\ref{sec:scope}). There is no other client operation.
In Lean, \lean{Translation.environment} fixes these data and
\lean{Dex.Environment.IsAdmissibleWord} states the shape. The \Lang{} form
is \lean{L.IsAdmissibleWord}: one \lean{start}, then \lean{input} requests.
The lemmas \lean{isAdmissibleWord_encode} and \lean{isAdmissibleWord_decode}
show that the encoded admissible \Lang{} words are exactly the words of
\Eadm.

Counting attempted requests, and not only accepted ones, matters. A
publication after the Flow closed is rejected and produces only a
$\ev{rejected}$ event; \Lang{} produces $\ev{refused}$ for the same input.
A word of accepted events only would give the \Lang{} run that exits and
then refuses an input no matching Dex run.

\paragraph{$\pi$.}
The projection $\pi : \Obs \to \ObsL \cup \{\varepsilon\}$
(\lean{Translation.project}) is given by Table~\ref{tab:projection}.

\begin{table}[t]
\centering
\small
\begin{tabular}{@{}ll@{}}
\toprule
Dex event & \Lang{} event \\
\midrule
$\ev{start}(T_0, \mathit{enc}(v))$ & $\ev{begin}(v)$ \\
$\ev{pub}(c, \mathit{enc}(v))$, $c$ external & $\ev{in}(c, v)$ \\
$\ev{rejected}(\ev{pub}(c, \mathit{enc}(v)))$, $c$ external & $\ev{refused}(c, v)$ \\
$\ev{fx}(e, m, a)$ & $\ev{out}(a)$ \quad (the source is dropped) \\
$\ev{closed}(\mathit{Completed}, [(e, \{\texttt{exit}: v\})])$ & $\ev{exit}(v)$ \\
$\ev{closed}(\mathit{Failed}(s), \_)$ & $\ev{fail}(s)$ \\
every other event & $\varepsilon$ \\
\bottomrule
\end{tabular}
\caption{The projection $\pi$ from Dex events to \Lang{} events.}
\label{tab:projection}
\end{table}

\subsection{Safety}

A finite trace of a system is the event sequence of a finite run from the
initial state, viewed through a function that may erase events
(\lean{TransitionSystem.IsFiniteTrace}). The two sets of
Equation~\eqref{eq:safety} are (\file{TraceSets.lean}):

\begin{Verbatim}
def IsProjectedDexFiniteTrace (program : Program) (trace : List Event) : Prop :=
  (Dex.system ⟦program⟧).IsFiniteTrace environment.IsAdmissiblePrefix project trace

def IsLFiniteTrace (program : Program) (trace : List Event) : Prop :=
  (L.system program).IsFiniteTrace (fun _ => True) some trace
\end{Verbatim}

Lean's core library has no set type, so an equality of sets is written as a
bi-implication for every element. The safety theorem
(\file{Theorems.lean}) is:

\begin{Verbatim}
theorem safety_of_isWellFormed (program : Program) (isWellFormed : program.IsWellFormed)
    (trace : List Event) :
    IsProjectedDexFiniteTrace program trace ↔ IsLFiniteTrace program trace
\end{Verbatim}

This is weak trace equality: Dex transitions without a projected event are
invisible. Finite traces alone cannot tell a Flow whose last thread ended
with \kw{DeadEnd}, and which stays running forever, from a Flow whose close
an unfair scheduler never runs. Liveness closes that gap.

\subsection{Liveness}
\label{sec:liveness-statement}

\paragraph{Runs for a word.}
Clock ticks are always enabled, so no maximal run is finite. A run
\emph{for} an input word $w$ is an infinite run whose request transitions
are exactly the requests of $w$, each once and in order
(\lean{IsRunFor}). A tick is time passing, not progress; every other
transition is a progress step.

\paragraph{Fairness.}
The rules themselves are unfair: the dispatcher may delay a startable
method, the timer service a due timer, and a Worker its response, each
without bound. Liveness is therefore stated for fair runs. The conditions
are assumptions about the scheduler, the timer service, and the Workers.
They are not guarantees that the rules provide. Table~\ref{tab:fairness}
states them on both sides. Each of F2 to F5 is a weak-fairness condition of
one shape (\lean{EventuallyFails}): a given predicate on the state does not
hold forever from any point on.

\begin{table}[t]
\centering
\small
\begin{tabularx}{\textwidth}{@{}P{0.6cm}LL@{}}
\toprule
 & Dex (\lean{Dex.IsFair}, \file{Dex/Observation.lean}) &
   \Lang{} (\lean{L.IsFair}, \file{L/Semantics.lean}) \\
\midrule
F0 & Every request of $w$ is eventually attempted, in order. It follows
  from ``run for $w$''. & The same. \\
F1 & The clock exceeds every bound. & The same. \\
F2 & No timer stays armed and due forever. An assumption on the timer
  service: it eventually proposes \rulename{TIMER-FIRE}. & No thread's
  armed timer stays due forever: \kw{fire} eventually happens. \\
F3 & No invocation stays pending forever. An assumption on dispatch: a
  pending method is eventually handed to a Worker. In this fragment every
  pending invocation is startable, because there are no locks. & No thread
  stays forever in a phase in which its next action has not begun
  (\lean{bootPending}, \lean{ready}, \lean{woken}). \\
F4 & No task stays in $\Theta$ forever without having returned: it emits
  all its effects and returns, or, if its invocation was canceled, it is
  dropped after some prefix of its effects. An assumption on the Workers:
  handlers terminate, and a Worker whose task was canceled eventually stops
  or finishes. & No thread stays \lean{emitting} forever while the program
  runs; after the end, a non-empty list of late emits eventually
  shrinks. \\
F5 & No task stays \emph{Returned} forever: its result commits, or it is
  dropped as stale. An assumption on the engine: a returned response is
  eventually proposed as a commit. & No thread stays forever in a phase in
  which it began an action and has not completed it. \\
F6 & Cross-Flow messages are delivered. Vacuous here: one Flow. &
  Vacuous. \\
\bottomrule
\end{tabularx}
\caption{The fairness assumptions. They constrain runs; no rule enforces
them.}
\label{tab:fairness}
\end{table}

F4 deserves emphasis. The server does not make a Worker finish. It hands a
method to a Worker, and for a canceled method it sends a cancel request;
cancellation is cooperative. That every task eventually leaves $\Theta$ is
an assumption about Workers that run terminating handlers and honor
cancellation or run to completion. A Worker that hangs forever on a task
gives an unfair run, about which the theorem says nothing.

\paragraph{Markers.}
A fair run is infinite, because time keeps passing. It is classified by its
progress steps (\lean{HasMarker}):
\begin{itemize}
\item \mterm{} (\lean{terminated}): finitely many progress steps, and the
  program is closed at some point.
\item \mdead{} (\lean{deadlocked}): finitely many progress steps, and the
  program is never closed. From some point on only ticks happen: nothing
  can start, no task is in flight, no timer is armed, and the word is used
  up. A program whose threads all ended without \kw{exit} lands here.
\item \mdiv{} (\lean{divergent}): infinitely many progress steps.
\end{itemize}
A \mterm{} run has finitely many progress steps, so by F4 and F5 no task
is left in $\Theta$ after the last one. The trace therefore already contains
every effect that a canceled Worker emits after the close, and a \mterm{}
trace is never a proper prefix of another fair trace for the same word.

\paragraph{The statement.}
A fair maximal trace for $w$ is the whole observable trace of a fair run
for $w$, with the run's marker (\lean{IsFairMaximalTrace}). The two sets of
Equation~\eqref{eq:liveness} are:

\begin{Verbatim}
def IsProjectedDexFairMaximalTrace (program : Program) (word : Sequence L.Request)
    (trace : Sequence Event) (marker : Marker) : Prop :=
  (Dex.system ⟦program⟧).IsFairMaximalTrace Dex.IsFair project
    (word.map encodeRequest) trace marker

def IsLFairMaximalTrace (program : Program) (word : Sequence L.Request)
    (trace : Sequence Event) (marker : Marker) : Prop :=
  (L.system program).IsFairMaximalTrace L.IsFair some word trace marker
\end{Verbatim}

A \lean{Sequence} is a finite or infinite sequence. The word is quantified
in \Lang{} form and encoded for the Dex side. The liveness theorem is:

\begin{Verbatim}
theorem liveness_of_isWellFormed (program : Program) (isWellFormed : program.IsWellFormed)
    (word : Sequence L.Request) (isAdmissible : L.IsAdmissibleWord word)
    (trace : Sequence Event) (marker : Marker) :
    IsProjectedDexFairMaximalTrace program word trace marker ↔
      IsLFairMaximalTrace program word trace marker
\end{Verbatim}

\lean{expressiveness_of_isWellFormed} is the conjunction of the two
theorems. The theorem names carry the scope, so that it is visible wherever
a theorem is used.

\paragraph{Why this equivalence.}
The specification chose trace equality over weak bisimilarity, which is
stronger than the claim needs and distinguishes \emph{where} a
nondeterministic choice is made. Trace equality with markers, per input
word, still separates normal exit, failure, deadlock, and divergence, and
it accounts for effects after the close. The proof does establish step
simulations in both directions over one relation (\S\ref{sec:proof}); the
theorems state only what follows for traces. The theorem is untimed: clock values are not observable. Time still
matters, through the order in which timers fire and through F1 and F2.

\paragraph{What the kernel checked.}
For both theorems, \texttt{\#print axioms} reports
\lean{propext}, \lean{Classical.choice}, and \lean{Quot.sound}: the three
standard axioms of Lean, and nothing else. No file contains \texttt{sorry},
and the proofs do not use \lean{native_decide}. The trusted base is the
Lean~4 kernel \cite{demoura2021lean4} and the definitions a reader must
check against the specification: the transition-system core, the Dex model,
the \Lang{} model, the translation, the projection, and the two statements,
about 3,450 lines in 14 files.
